\textbf{Phần 1. Trắc nghiệm lựa chọn}
\setcounterpageref{ex}{0}

\begin{ex} %[4D12V64]
Phân số đảo ngược của phân số $\dfrac{3}{5}$ là
\choice
{$\dfrac{3}{5}$}
{\True $\dfrac{5}{3}$}
{$-\dfrac{3}{5}$}
{$\dfrac{1}{5}$}
\loigiai{
Đảo ngược tử số và mẫu số: $\dfrac{5}{3}$.\\
\textbf{Đáp án của câu hỏi là B}
}
\end{ex}

\begin{ex} %[4D12V64]
Muốn chia một phân số cho một phân số, ta
\choice
{\True Lấy phân số thứ nhất nhân với phân số thứ hai đảo ngược}
{Lấy tử số chia tử số, mẫu số chia mẫu số}
{Lấy mẫu số nhân tử số}
{Chia hai tử số với nhau}
\loigiai{
Nhân phân số thứ nhất với phân số thứ hai đảo ngược.\\
\textbf{Đáp án của câu hỏi là A}
}
\end{ex}

\begin{ex} %[4D12V64]
Kết quả của phép tính $\dfrac{2}{3} : \dfrac{3}{4}$ là
\choice
{$\dfrac{6}{12}$}
{$\dfrac{1}{2}$}
{\True $\dfrac{8}{9}$}
{$\dfrac{9}{8}$}
\loigiai{
$\dfrac{2}{3} \times \dfrac{4}{3} = \dfrac{8}{9}$.\\
\textbf{Đáp án của câu hỏi là C}
}
\end{ex}

\begin{ex} %[4D12V64]
Kết quả của phép tính $\dfrac{4}{5} : 2$ là
\choice
{$\dfrac{8}{5}$}
{$\dfrac{4}{10}$}
{$\dfrac{1}{5}$}
{\True $\dfrac{2}{5}$}
\loigiai{
$\dfrac{4}{5} \times \dfrac{1}{2} = \dfrac{4}{10} = \dfrac{2}{5}$.\\
\textbf{Đáp án của câu hỏi là D}
}
\end{ex}

\begin{ex} %[4D12V64]
Kết quả của phép tính $3 : \dfrac{1}{2}$ là
\choice
{$\dfrac{3}{2}$}
{\True $6$}
{$\dfrac{1}{6}$}
{$1$}
\loigiai{
$3 \times \dfrac{2}{1} = 6$.\\
\textbf{Đáp án của câu hỏi là B}
}
\end{ex}

\begin{ex} %[4D12V64]
Tìm $x$ biết: $x \times \dfrac{2}{5} = \dfrac{4}{7}$.
\choice
{\True $\dfrac{10}{7}$}
{$\dfrac{8}{35}$}
{$\dfrac{7}{10}$}
{$\dfrac{6}{12}$}
\loigiai{
$x = \dfrac{4}{7} : \dfrac{2}{5} = \dfrac{4}{7} \times \dfrac{5}{2} = \dfrac{10}{7}$.\\
\textbf{Đáp án của câu hỏi là A}
}
\end{ex}

\begin{ex} %[4D12V64]
Một hình chữ nhật có diện tích $\dfrac{2}{3}\text{ m}^2$, chiều dài $\dfrac{4}{5}\text{ m}$. Chiều rộng của hình chữ nhật là
\choice
{$\dfrac{8}{15}\text{ m}$}
{$\dfrac{6}{5}\text{ m}$}
{\True $\dfrac{5}{6}\text{ m}$}
{$\dfrac{1}{3}\text{ m}$}
\loigiai{
$\dfrac{2}{3} : \dfrac{4}{5} = \dfrac{2}{3} \times \dfrac{5}{4} = \dfrac{5}{6}\text{ m}$.\\
\textbf{Đáp án của câu hỏi là C}
}
\end{ex}

\begin{ex} %[4D12V64]
Có $\dfrac{7}{2}\text{ l}$ dầu chia đều vào các chai, mỗi chai chứa $\dfrac{1}{2}\text{ l}$ dầu. Số chai dầu rót được là
\choice
{$6$ chai}
{$8$ chai}
{$9$ chai}
{\True $7$ chai}
\loigiai{
$\dfrac{7}{2} : \dfrac{1}{2} = \dfrac{7}{2} \times 2 = 7$ chai.\\
\textbf{Đáp án của câu hỏi là D}
}
\end{ex}

\begin{ex} %[4D12V64]
Một tấm vải dài $\dfrac{8}{3}\text{ m}$ được cắt thành các đoạn dài $\dfrac{2}{3}\text{ m}$. Số đoạn vải cắt được là
\choice
{$3$ đoạn}
{\True $4$ đoạn}
{$5$ đoạn}
{$6$ đoạn}
\loigiai{
$\dfrac{8}{3} : \dfrac{2}{3} = 4$ đoạn.\\
\textbf{Đáp án của câu hỏi là B}
}
\end{ex}

\begin{ex} %[4D12V64]
Tìm $x$ biết: $x : \dfrac{3}{8} = \dfrac{4}{9}$.
\choice
{\True $\dfrac{1}{6}$}
{$\dfrac{32}{27}$}
{$\dfrac{7}{17}$}
{$\dfrac{12}{72}$}
\loigiai{
$x = \dfrac{4}{9} \times \dfrac{3}{8} = \dfrac{12}{72} = \dfrac{1}{6}$.\\
\textbf{Đáp án của câu hỏi là A}
}
\end{ex}

\begin{ex} %[4D12V64]
Tính giá trị biểu thức: $\dfrac{3}{5} : \dfrac{2}{3} \times \dfrac{4}{9}$.
\choice
{$\dfrac{1}{5}$}
{$\dfrac{3}{5}$}
{\True $\dfrac{2}{5}$}
{$\dfrac{4}{5}$}
\loigiai{
$\dfrac{3}{5} \times \dfrac{3}{2} \times \dfrac{4}{9} = \dfrac{9}{10} \times \dfrac{4}{9} = \dfrac{4}{10} = \dfrac{2}{5}$.\\
\textbf{Đáp án của câu hỏi là C}
}
\end{ex}

\begin{ex} %[4D12V64]
Một hình bình hành có diện tích $\dfrac{5}{8}\text{ m}^2$, chiều cao $\dfrac{3}{4}\text{ m}$. Độ dài đáy của hình bình hành là
\choice
{$\dfrac{15}{32}\text{ m}$}
{$\dfrac{6}{5}\text{ m}$}
{$\dfrac{1}{2}\text{ m}$}
{\True $\dfrac{5}{6}\text{ m}$}
\loigiai{
$a = S : h = \dfrac{5}{8} : \dfrac{3}{4} = \dfrac{5}{8} \times \dfrac{4}{3} = \dfrac{5}{6}\text{ m}$.\\
\textbf{Đáp án của câu hỏi là D}
}
\end{ex}

\textbf{Phần 2. Trắc nghiệm đúng sai}
\setcounterpageref{ex}{0}

\begin{ex} %[4D12V64]
Xét tính đúng sai của các phép chia phân số:
\choiceTF
{\True Phân số đảo ngược của $\dfrac{4}{7}$ là $\dfrac{7}{4}$.}
{\True $\dfrac{3}{5} : \dfrac{3}{5} = 1$.}
{$\dfrac{2}{3} : 2 = \dfrac{4}{3}$.}
{\True $\dfrac{2}{3} : 2 = \dfrac{1}{3}$.}
\loigiai{
\begin{itemize}
    \item Đảo ngược đúng.
    \item Chia cho chính nó bằng 1: Đúng.
    \item $2/3 : 2 = 1/3 \neq 4/3$. Ý 3 Sai, Ý 4 Đúng.
\end{itemize}
\begin{itemize}
    \item a) \textbf{Đúng}.
    \item b) \textbf{Đúng}.
    \item c) \textbf{Sai}.
    \item d) \textbf{Đúng}.
\end{itemize}
}
\end{ex}

\begin{ex} %[4D12V64]
Người ta chia một sợi dây dài $\dfrac{15}{4}\text{ m}$ thành các đoạn nhỏ bằng nhau:
\choiceTF
{\True Nếu chia thành 3 đoạn bằng nhau thì mỗi đoạn dài $\dfrac{5}{4}\text{ m}$.}
{\True Nếu chia thành các đoạn dài $\dfrac{3}{4}\text{ m}$ thì được 5 đoạn.}
{Nếu chia đôi thì mỗi đoạn dài $\dfrac{15}{2}\text{ m}$.}
{\True Nếu chia đôi thì mỗi đoạn dài $\dfrac{15}{8}\text{ m}$.}
\loigiai{
\begin{itemize}
    \item $15/4 : 3 = 5/4\text{ m}$. Đúng.
    \item $15/4 : 3/4 = 5$ đoạn. Đúng.
    \item Chia đôi: $15/4 : 2 = 15/8\text{ m} \neq 15/2\text{ m}$. Ý 3 Sai, Ý 4 Đúng.
\end{itemize}
\begin{itemize}
    \item a) \textbf{Đúng}.
    \item b) \textbf{Đúng}.
    \item c) \textbf{Sai}.
    \item d) \textbf{Đúng}.
\end{itemize}
}
\end{ex}

\textbf{Phần 3. Trả lời ngắn}
\setcounterpageref{ex}{0}

\begin{ex} %[4D12V64]
Tính thương: $\dfrac{3}{7} : \dfrac{5}{7}$ (dạng a/b).
\shortans{3/5}
\loigiai{
\textbf{Đáp số:} $3/5$\\
$3/5$.
}
\end{ex}

\begin{ex} %[4D12V64]
Tính thương: $\dfrac{4}{9} : \dfrac{2}{3}$ (rút gọn tối giản dạng a/b).
\shortans{2/3}
\loigiai{
\textbf{Đáp số:} $2/3$\\
$4/9 \times 3/2 = 2/3$.
}
\end{ex}

\begin{ex} %[4D12V64]
Tìm $x$ biết: $x \times \dfrac{3}{4} = \dfrac{9}{16}$ (dạng a/b).
\shortans{3/4}
\loigiai{
\textbf{Đáp số:} $3/4$\\
$x = 9/16 : 3/4 = 3/4$.
}
\end{ex}

\begin{ex} %[4D12V64]
Một hình chữ nhật có diện tích $\dfrac{3}{5}\text{ m}^2$, chiều rộng $\dfrac{1}{2}\text{ m}$. Tính chiều dài (dạng a/b).
\shortans{6/5}
\loigiai{
\textbf{Đáp số:} $6/5$\\
$3/5 : 1/2 = 6/5\text{ m}$.
}
\end{ex}

\begin{ex} %[4D12V64]
Có $6\text{ kg}$ đường chia đều vào các túi $\dfrac{3}{4}\text{ kg}$. Tính số túi đường đóng được.
\shortans{8}
\loigiai{
\textbf{Đáp số:} $8$\\
$6 : (3/4) = 8$ túi.
}
\end{ex}

\begin{ex} %[4D12V64]
Tính giá trị biểu thức: $\dfrac{1}{2} : \dfrac{2}{3} : \dfrac{3}{4} : \dfrac{4}{5}$ (dạng a/b).
\shortans{5/4}
\loigiai{
\textbf{Đáp số:} $5/4$\\
$1/2 \times 3/2 \times 4/3 \times 5/4 = 5/4$.
}
\end{ex}
